\chapter{可微流形}\label{chap:differentiable-manifolds}
\begingroup
\renewcommand{\thefootnote}{}
\footnotetext{\textcopyright\ Springer International Publishing Switzerland 2014.
  L. Godinho and J. Natário, \emph{An Introduction to Riemannian Geometry},
  Universitext. DOI \href{https://doi.org/10.1007/978-3-319-08666-8_1}{10.1007/978-3-319-08666-8\_1}.}
\endgroup

在纯粹数学和应用数学中, 我们经常遇到局部看起来像 $\mathbb{R}^n$ 的空间, 意思是这些空间在局部可以用 $n$ 个坐标来参数化. 例如, $n$ 维球面 $S^n\subset\mathbb{R}^{n+1}$, 或刚体的构型集合 $\mathbb{R}^3\times SO(3)$. 我们或许会预期微积分的基本工具在这些空间上仍然适用. 然而, 由于一般没有典范的局部坐标选取, 在讨论导数或积分等概念时必须格外小心: 这些概念在 $\mathbb{R}^n$ 中的定义依赖于优先选取的 Cartesian 坐标.

本章的主题是这类空间 (称为\textbf{可微流形}) 的严格定义, 以及与之有关的微分概念. 尽管直观想法看来相当简单, 而且事实上可以追溯到 Gauss 和 Riemann, 但直到 1936 年才由 Whitney 给出形式定义.

局部看起来像 $\mathbb{R}^n$ 的空间, 通过\textbf{拓扑流形}的定义加以形式化: 拓扑流形是局部同胚于 $\mathbb{R}^n$ 的拓扑空间. 我们将在第~\RGref{sec:topological-manifolds}~节研究它们, 并讨论若干例子, 尤其是 2 维的例子 (曲面).

第~\RGref{sec:differentiable-manifolds}~节把可微流形定义为坐标变换 (从 $\mathbb{R}^n$ 到 $\mathbb{R}^n$ 的映射) 光滑 ($C^\infty$) 的拓扑流形. 由此可以把\textbf{可微函数}定义为局部坐标表达式光滑的函数 (第~\RGref{sec:differentiable-maps}~节), 把\textbf{切向量}定义为作用于实值可微函数的方向导数算子 (第~\RGref{sec:tangent-space}~节). 第~\RGref{sec:immersions-embeddings}~节考察可微映射的两个重要例子, 即\textbf{浸入}和\textbf{嵌入}.

\textbf{向量场}及其\textbf{流}是第~\RGref{sec:vector-fields}~节的主要内容. 我们定义向量场之间的一种自然微分运算, 称为 \textbf{Lie 括号}; 它度量这些向量场的流的不交换性, 在微分几何中起着核心作用.

第~\RGref{sec:lie-groups}~节讨论一类重要的可微流形, 它们同时也是群, 称为 \textbf{Lie 群}. 我们将说明, 每个 Lie 群都可对应一个 \textbf{Lie 代数}, 即一个带有 Lie 括号的向量空间. 本节也讨论在 Lie 群作用下得到的流形的\textbf{商}.

流形的\textbf{可定向性} (与曲面"有两面"这一直观概念密切相关) 和\textbf{带边界流形} (推广有闭曲线作为边界的曲面, 或有闭曲面作为边界的立体这一概念), 分别在第~\RGref{sec:orientability}~节和第~\RGref{sec:manifolds-with-boundary}~节研究. 陈述著名的 Stokes 定理需要这两个概念, 该定理将在第~\RGref{chap:differential-forms}~章证明.
